
Growing neural networks is an emerging research direction, offering potential benefits in training time reduction, improved optimization, and a frugal form of task-aware architecture adaptation. As shown in Sec.~\ref{sec:results}, growth can yield substantial reductions in training time, especially when growing depth, and these results translate to transformer architectures as well (Sec.~\ref{sec:results_transformers}). An important next step is to characterize the settings in which growth is most effective relative to complementary approaches such as pruning. In particular, systematic comparisons on modern GPU hardware would clarify how their respective computational advantages translate into runtime and energy savings~\citep{boumendil_grow_2023}, and how these methods can be combined.

Existing methods optimize a local proxy objective, such as selecting the steepest immediate loss decrease, following the natural gradient, or heuristics intended to preserve trainability. An open question is how closely the proxy objective optimized by repeated growing operations of Eq.~\eqref{eqn:grow_decomposition} approximates the idealized objective of Eq.~\eqref{eqn:ideal_obj}. Treating growth as a sequential decision-making problem, one can decide where, when, and how to grow to more directly optimize Eq.~\eqref{eqn:ideal_obj}. This is similar to Reinforcement Learning approaches used in neural architecture search~\citep{zoph_neural_2017}.

Optimizer-aware growth methods are another promising direction: it is common practice to simply reset the optimizer state upon growth. Variance Transfer~\citep{yuan_accelerated_2023} is a step in the right direction, adapting the learning rate (and time dependence of adaptive optimizers, such as Adam) to the growth stage. Future work could develop principled strategies for initializing the optimizer state (\eg{} momentum) of newly grown parameters, beyond zero initialization or, when duplicating parameters, copying the state from the source parameters~\citep{shen_staged_2022}.

As in many research domains, a large-scale benchmark would benefit the whole community. Open, centralized, and easy-to-reproduce benchmarks would make it easier to compare new ideas with the existing literature. Indeed, an overview of existing methods is the first step, and this is the purpose of this survey. We hope that this prompts the next step: a reference benchmark, open to contributions, with appropriate metrics that encompass the multiple objectives relevant to growth.

Such benchmarks would be especially useful for the most ambitious goal of growth: frugal architecture adaptation. Existing hand-crafted architectures reflect decades of research, and are already highly-tuned to the standard datasets considered here. A fairer (and more realistic) test would evaluate growth on ``unseen'' datasets, similar to existing benchmarks for NAS~\citep{geada_unseennaschallenge_2024}. As the optimal architecture is not known in advance and may lie far from the initial seed architecture, this would require principled criteria for \emph{where} to grow and \emph{when} to stop, which are relatively understudied. We hope that the ideas in this survey will spur the community to grow in promising new directions.
